\begin{tikzpicture}[font=\figurefont\small,x=8mm,y=8mm,
  route/.style={draw=BookTeal,line width=1.5pt,-{Stealth[length=2mm]}},
  high/.style={draw=BookInk,fill=white,minimum size=5mm,inner sep=1pt}]
  \draw[step=1,BookRule!70] (.5,.5) grid (8.5,5.5);
  \draw[BookInk,line width=2pt] (4.5,.5) -- (4.5,1.65);
  \draw[BookInk,line width=2pt] (4.5,2.35) -- (4.5,3.65);
  \draw[BookInk,line width=2pt] (4.5,4.35) -- (4.5,5.5);
  \draw[route] (2,2) -- (4,2);
  \draw[route] (4,2) -- (7,2);
  \foreach \x in {2,3,4,5,6,7} \fill[BookTeal] (\x,2) circle (1.6pt);
  \draw[BookAmber,dashed,line width=1.2pt,-{Stealth[length=2mm]}]
    (2,2) -- (2,4) -- (7,4) -- (7,2);
  \node[high] at (2,2) {S};
  \node[high] at (4,2) {D};
  \node[high,fill=FigureOutput!12] at (7,2) {G};
  \node[above,fill=white,inner sep=1pt] at (4,4) {强制子目标};
  \node at (2.5,6.1) {左房间};
  \node at (6.5,6.1) {右房间};
  \node[align=left,anchor=west,text width=65mm] at (9.7,4.5)
    {实线：五次原子动作\\高层：先到D，再到G\\下门位于$y=2$};
  \node[align=left,anchor=west,text width=65mm] at (9.7,2)
    {虚线：先到$(4,4)$\\最短也需九步\\上门位于$y=4$};
  \node[align=center,text width=145mm] at (9,-.8)
    {方框S、D处选择高层行为；中间圆点处仍由内部策略选择原子动作。};
\end{tikzpicture}
